% !TEX root = ../main.tex


\chapter{课题的来源及研究目的和意义}

\section{课题来源}

在现代存储系统中，数据可用性对于应对各种故障至关重要\cite{fordAvailabilityGloballyDistributed2010}。
目前，主流的容错技术包括数据复制和纠删码。
数据复制通过在不同节点保存多份副本，虽然其实现简单，但是存储开销较大。
例如，三副本机制能同时容忍小于等于两个故障，但是需要三倍的存储空间。
相比之下，纠删码技术通过对数据进行分块编码和解码，能够以远低于复制的存储成本，提供同等甚至更高的数据可靠性\cite{weatherspoonErasureCodingVs2002}。
现如今，纠删码已成为Google文件系统（GFS）\cite{ghemawatGoogleFileSystem2003}、Windows Azure存储\cite{calderWindowsAzureStorage2011, huangErasureCodingWindows2012}、Hadoop分布式文件系统（HDFS）\cite{shvachkoHadoopDistributedFile2010, chiniahDynamicErasureCoding2019}以及Ceph\cite{weilCephScalableHighPerformance2006, weilCephReliableScalable2007}等大规模存储系统的重要技术。

纠删码通过对原始数据分块进行编码生成冗余校验信息，从而提升数据可靠性，并在数据丢失时实现恢复。
与复制技术相比，经典纠删码Reed-Solomon码（RS码）\cite{reedPolynomialCodesCertain1960}的编码和解码过程计算量更大\cite{weatherspoonErasureCodingVs2002}，其通常依赖于有限域$GF[2^w], w>1$上的矩阵乘法，并在计算时将部分有限域上的运算转换为异或运算（亦即有限域$GF[2]$上的运算）\cite{lubyXORBasedErasureResilientCoding1995}从而提高计算速度。
众多研究如Jerasure\cite{plankJerasureLibraryFacilitating2008, arslanOpenMPPOSIXThreads2017}、G-CRS\cite{liuGCRSGPUAccelerated2018}、Zerasure\cite{zhouFastErasureCoding2019, zhouFastErasureCoding2020}、SLPEC\cite{uezatoAcceleratingXORbasedErasure2021}、Cerasure\cite{niuCerasureFastAcceleration2023, wangFastAccelerationStrategies2025}和Intel的ISA-L库\cite{gajalwarExploringDataProtection2024}等都致力于提升基于有限域矩阵乘法的编解码效率，以降低系统成本并增强服务可靠性。
近年来，基于非循环移位异或运算的纠删码在理论层面取得许多进展，例如ZigZag Decodable Codes\cite{sungZigZagDecodableCodeMDS2013, daiNewZigzagDecodableCode2017, gongZigzagDecodableCodes2017}和Two-tone Shift-XOR Codes\cite{fuOverheadfreeInplaceRecovery2014, fuOverheadfreeInplaceRecovery2015, fuDecodingRepairSchemes2020, chengSuccessivelySolvableShiftAdd2021, fuTwotoneShiftXORStorage2021}，并且在存储\cite{gongZigzagDecodableCodes2017, daiBandwidthOverheadFreeData2017}、密码学\cite{gongZigzagDecodableRampSecret2018, kabiriradEfficientRampSecret2023, chenNewShiftAddSecret2024}和网络\cite{daiZigzagDivisionMultipleAccess2019}等应用领域进行了各种拓展研究。
对于这类移位异或的编解码算法来说，异或运算就是最基本的运算，而并非像传统纠删码那样需要转化，这种更简单更直接的运算方式也取得了更高的理论性能\cite{gongZigzagDecodableCodes2017}。
相比于ZigZag算法，Two-tone算法对于移位和异或操作的编排更具有规则性，使得其在解码时无需多次扫描即可解出数据的位置，更容易按顺序访问内存，在算法实现层面有更高效的潜能\cite{chengSuccessivelySolvableShiftAdd2021, fuTwotoneShiftXORStorage2021}。

尽管Two-tone移位异或纠删码算法在理论上具备优越性能，其完整的工程实现与系统级优化方面的工作尚且缺少。现有文献虽展示了算法潜力，但缺乏可用于生产环境的提速方案和高性能实现。具体来讲，可总结为以下三点问题：

\begin{itemize}
    \item 缺少泛化性实现：当前的算法实现只能在实验测试中证明高效性，但并不具有泛化能力，无法应对数据块和校验块（存储冗余信息的块）所有的丢失情况，不能同时满足存储系统中解码丢失数据块和重新编码丢失校验块的双重需求。
    \item 缺少系统级优化：理论算法没有明确给出程序优化的策略，例如减少条件分支和遍历次数、增强缓存局部性和减轻内存访问压力等方面的优化，同时RS纠删码的系统级优化策略因其基本算法不同而无法直接应用到移位异或纠删码上。
    \item 缺少工业级应用：上述纠删码研究成果仅有Jerasure和ISA-L等少数基于RS码的算法以插件的形式集成到了Ceph和Hadoop等开源存储系统中，大部分纠删码（包括移位异或纠删码）目前并没有在工业级应用上测试并验证高效性。
\end{itemize}

鉴于上述三点问题，本课题将Two-tone非循环移位异或纠删码以插件形式集成至开源分布式存储系统Ceph中，作为其存储池管理的一种存储编码，借助Ceph中的应用更加规范地对比测试纠删码算法，并在该环境下应用各种系统级优化方案，提炼出针对移位异或纠删码的高效实现方法论，验证各方法论的有效性，达成比Ceph中的SOTA纠删码插件和最新研究中的SOTA纠删码算法更高效的性能。

\section{研究的目的和意义}

本课题旨在设计高效的Two-tone非循环移位异或纠删码实现方案，推动这类新型移位异或纠删码算法在大规模分布式存储系统例如Ceph中的落地应用，实现比其他纠删码算法更加优秀的应用性能。
通过理论算法与工程实践的创新性结合，力求进一步提高纠删码在实际应用中的性能，让大规模分布式存储系统更加高效。

本课题的研究目的一是针对Two-tone移位异或纠删码算法本身在程序实现时的瓶颈，提出系统级优化方法，包括减少条件分支和循环遍历次数、提升缓存命中率、优化内存访问模式等，提炼并总结适用于新型移位异或纠删码的优化方案，为后续更复杂深入的优化打下基础；
二是将Two-tone纠删码以插件的形式集成至Ceph分布式存储系统中，通过完整的功能性测试，验证各优化方案的有效性，并实施性能对比实验，与Ceph中的SOTA纠删码ISA-L比较，与RS码的SOTA实现Cerasure和移位异或码的SOTA实现ZigZag比较，取得更高效的性能。

本课题的研究意义一是推动非循环移位异或纠删码从算法理论到实际应用的发展，呈现新型高效纠删码在存储系统中的解决方案，丰富存储系统的编码手段，并实现更高效的纠删码算法，进一步提高编解码吞吐量，促进解决纠删码吞吐量不匹配高性能设备的问题\cite{weiCharacterizingOptimizingRemote2021}，为分布式存储系统的性能提升提供技术支撑；
二是填补移位异或纠删码缺少系统级优化方法论的空白，对齐RS码的研究进度，发挥出移位异或码的高效潜能，同时为开源存储系统Ceph提供此前未有的高性能移位异或纠删码插件，促进开源社区的发展，推动相关领域的技术进步。

% 本课题的研究目的主要在于以下三点：

% \begin{itemize}
%     \item 针对Two-tone移位异或纠删码算法本身在程序实现时的瓶颈，提出系统级优化方法，包括减少条件分支和循环次数、提升缓存命中率、优化内存访问模式等，提炼并总结适用于新型移位异或纠删码的优化方案。
%     \item 将Two-tone纠删码以插件的形式集成至Ceph分布式存储系统中，通过完整的功能性测试，验证各优化方案的有效性，并部署至存储池加以应用测试。
%     \item 实施性能对比实验，与Ceph中的SOTA纠删码ISA-L比较，与RS码的SOTA实现Cerasure和移位异或码的SOTA实现ZigZag比较，取得更高效的性能。
% \end{itemize}

% 本课题的研究意义主要在于以下四点：

% \begin{itemize}
%     \item 实现更高效的纠删码算法，进一步提高编解码吞吐量，促进解决纠删码吞吐量不匹配高性能设备的问题\cite{weiCharacterizingOptimizingRemote2021}，为分布式存储系统的性能提升提供技术支撑。
%     \item 推动非循环移位异或纠删码从算法理论到实际应用，呈现新型高效纠删码在存储系统中的解决方案，丰富存储系统的编码手段，拓展移位异或纠删码的应用场景。
%     \item 填补移位异或纠删码缺少系统级优化方法论的空白，对齐RS码的研究进度，发挥出移位异或码的高效潜能，推动相关领域的技术进步。
%     \item 为开源存储系统Ceph提供可扩展、高性能的纠删码插件，促进开源社区的发展，推动国产分布式存储系统的自主创新和产业化应用。
% \end{itemize}

% 综上所述，本课题不仅具有重要的理论价值，更具备显著的工程应用前景。通过系统性研究和优化，能够为分布式存储系统提供更高效、更可靠的数据保护方案，推动新型纠删码技术的实际落地与广泛应用。



